% ----------------------------------------------------------------------
%  2.1.4. Образование и отжиг точечных дефектов
% ----------------------------------------------------------------------
\subsection{Образование и отжиг точечных дефектов}\label{ssec:annealing}

\begin{orientir}
Равновесная концентрация вакансий меняется с температурой. Но откуда
берутся новые вакансии при нагреве и куда деваются лишние при охлаждении?
Здесь мы разберём: механизмы рождения дефектов (Френкеля и Шоттки);
источники и стоки; как быстро кристалл приходит к равновесию; как по
экспериментальным кривым отжига определить энергию активации.
\end{orientir}

\subsubsection{Механизмы Френкеля и Шоттки}

Равновесие точечных дефектов в кристалле \emph{динамическое}: дефекты
непрерывно рождаются и исчезают, а их средняя концентрация постоянна.
Рождаться они могут двумя путями (рис.~\ref{fig:frenkel-schottky}).

\begin{opredelenie}
\term{Механизм Френкеля}: атом уходит из узла в междоузлие; одновременно
образуются вакансия и межузельный атом~— \term{пара Френкеля}.
\term{Механизм Шоттки}: атом с поверхности уходит наружу (испаряется или
переходит в адсорбированный слой), его место занимает атом из более
глубокого слоя, и т.\,д.~— вакансия «уходит» в глубь кристалла.
\end{opredelenie}

\begin{figure}[ht]
\centering
\begin{subfigure}[b]{0.36\textwidth}\centering
\begin{tikzpicture}[scale=0.55]
  \foreach \x in {0,...,5}{\foreach \y in {0,...,4}{
    \pgfmathtruncatemacro{\code}{10*\x+\y}
    \ifnum\code=22 \node[vac=4.6mm] at (\x,\y) {};
    \else \node[atom=4.6mm] at (\x,\y) {};\fi}}
  \node[atomI=4.6mm] at (3.5,2.5) {};
  \draw[arr, fmRed] (2.2,2.15) .. controls (2.7,2.6) .. (3.25,2.52);
\end{tikzpicture}
\caption{механизм Френкеля}
\end{subfigure}\hfill
\begin{subfigure}[b]{0.6\textwidth}\centering
\begin{tikzpicture}[scale=0.42]
  \foreach \s/\vx/\vy/\ad in {0/9/9/0, 7.5/3/4/1, 15/3/3/1}{
    \begin{scope}[shift={(\s,0)}]
    \fill[fmGray!12] (-0.6,-0.6) rectangle (5.6,4.5);
    \foreach \x in {0,...,5}{\foreach \y in {0,...,4}{
      \pgfmathtruncatemacro{\ok}{(\x==\vx && \y==\vy) ? 0 : 1}
      \ifnum\ok=1 \node[atom=3.5mm] at (\x,\y) {};\else\node[vac=3.5mm] at (\x,\y) {};\fi}}
    \ifnum\ad=1 \node[atomB=3.5mm] at (3,5.05) {};\fi
    \end{scope}}
  \node[atomB=3.5mm] at (3,4) {};
  \draw[arr, fmRed] (3,4.4) -- (3,5.4);
  \draw[arr, fmRed] (10.5,3.4) -- (10.5,3.85);
  \draw[arr, fmGray] (5.9,2) -- (6.7,2); \draw[arr, fmGray] (13.4,2) -- (14.2,2);
  \node[slbl] at (2.5,-1.3) {1}; \node[slbl] at (10,-1.3) {2}; \node[slbl] at (17.5,-1.3) {3};
\end{tikzpicture}
\caption{механизм Шоттки}
\end{subfigure}
\caption{Механизмы рождения точечных дефектов: \textit{а}~— пара Френкеля;
\textit{б}~— вакансия Шоттки: (1) поверхностный атом выходит в
адсорбированный слой, (2) на его место переходит атом из глубины,
(3) вакансия продвигается внутрь кристалла}
\label{fig:frenkel-schottky}
\end{figure}

Механизм Френкеля требует энергии $\EvF+\EiF\approx4$--$5\eV$, поэтому
в тепловом равновесии он практически не работает. Зато именно так
образуются дефекты при облучении частицами (п.~\ref{ssec:production}).
Механизм Шоттки гораздо легче: при выходе атома на поверхность все связи
не рвутся. Образно говоря, \emph{кристалл «растворяет» в себе пустоту}.

\subsubsection{Источники и стоки}

\begin{opredelenie}
\term{Источник} точечных дефектов~— место, где они рождаются;
\term{сток}~— место, куда они уходят и где исчезают.
\end{opredelenie}

Источники тепловых вакансий:
\begin{itemize}
\item свободная поверхность кристалла, а также внутренние поры и трещины;
\item границы зёрен (в поликристалле они работают так же эффективно, как
поверхность);
\item дислокации (вакансии рождаются и поглощаются при переползании,
п.~\ref{ssec:disl-motion}).
\end{itemize}
Если кристалл \emph{пересыщен} дефектами, те же места работают как стоки:
появляется термодинамическая движущая сила, которая снижает концентрацию
дефектов до равновесной.

\begin{opredelenie}
\term{Отжиг дефектов}~— процесс исчезновения избыточных дефектов из
пересыщенного ими кристалла.
\end{opredelenie}
Избыточные дефекты исчезают двумя путями:
\begin{enumerate}[label=\arabic*)]
\item \vazhno{уход на стоки}~— единственный путь, если в кристалле дефекты
одного типа (например, только вакансии и их комплексы);
\item \vazhno{рекомбинация}~— взаимное уничтожение дефектов разного типа,
прежде всего встреча вакансии и межузельного атома.
\end{enumerate}

\subsubsection{Как быстро кристалл приходит к равновесию}

Лишь стоки с дальнодействующим полем напряжений (например, дислокации)
притягивают вакансии издалека. В остальные стоки вакансия попадает,
блуждая случайно. На расстояниях, много больших межатомного, такое
блуждание описывается законами диффузии:
\begin{equation}\label{eq:Fick-v}
\vb J_v=-D_v\nabla C_v,\qquad
\frac{\partial C_v}{\partial t}=D_v\nabla^2C_v .
\end{equation}
Второе уравнение (\term{второй закон Фика}) по форме совпадает с
уравнением теплопроводности, для которого известно множество
аналитических решений. Обычно нужно знать не всё поле $C_v(\vb r,t)$, а
лишь среднюю по объёму избыточную концентрацию $\Delta\bar C_v(t)$. Если
сток одного типа и побочных реакций нет, она убывает экспоненциально:
\begin{keyeq}\label{eq:tau-sink}
\Delta\bar C_v(t)=\Delta\bar C_v(0)\,e^{-t/\tau},\qquad
\tau=\alpha\,\frac{r^2}{D_v},
\end{keyeq}
где $r$~— характерный размер «зоны питания» стока (половина среднего
расстояния между стоками), $\alpha$~— безразмерный множитель, зависящий
от геометрии (табл.~\ref{tab:alpha-sinks}).

\begin{smysl}
Формулу $\tau\sim r^2/D_v$ можно получить из одной размерности: время,
за которое частица диффузионно проходит расстояние $r$, порядка
$r^2/D$. Число стоков и их форма влияют лишь на множитель $\alpha$.
\end{smysl}

\begin{table}[ht]
\centering
\caption{Геометрический множитель $\alpha$ в формуле~\eqref{eq:tau-sink}
($r_0$~— радиус захвата стока)}
\label{tab:alpha-sinks}
\small
\begin{tabular}{@{}lll@{}}
\toprule
Сток & Зона питания & $\alpha$ \\
\midrule
Граница зерна (поверхность сферы радиусом $r$) & шар & $1/\pi^2$ \\
Свободная поверхность пластины толщиной $2r$ & плоский слой & $4/\pi^2$ \\
Точечный сток (пора, частица) & шар радиусом $r$ & $r/(3r_0)$ \\
Линейный сток (дислокация) & цилиндр радиусом $r$ & $\frac12\ln(r/r_0)$ \\
\bottomrule
\end{tabular}
\end{table}

Из~\eqref{eq:tau-sink} следуют важные выводы\ispr{В оригинале (с.~180)
для слоя толщиной $2r$ указано $\alpha=1/\pi^2$; для такой геометрии
наименьшая мода затухания $\cos(\pi x/2r)$ даёт $\alpha=4/\pi^2$.}.
\begin{itemize}
\item При нескольких типах стоков побеждает тот, у которого $\alpha r^2$
меньше,~— обычно ближайший.
\item $\tau$ зависит от температуры так же, как $D_v$:
$\tau\propto\exp(\Evm/kT)$. По скорости отжига можно найти $\Evm$.
\item $\tau$ \emph{не зависит} от начального пересыщения.
\end{itemize}

\begin{primer}[Время отжига вакансий в меди]
Пусть $\Evm=0{,}7\eV$, $a=2{,}5\cdot10^{-10}$~м, $\nu_0=10^{13}$~с$^{-1}$.
При $T=300$~К
$D_v=a^2\nu_0e^{-\Evm/kT}\approx6\cdot10^{-7}\cdot1{,}8\cdot10^{-12}
\approx1\cdot10^{-18}$~м$^2$/с.
\begin{itemize}
\item \emph{Стоки~— границы зёрен}, размер зерна 50~мкм ($r=25$~мкм):
$\tau\approx\frac{1}{\pi^2}\cdot\frac{(25\cdot10^{-6})^2}{10^{-18}}\approx
6\cdot10^{7}$~с~— около двух лет.
\item \emph{Стоки~— дислокации} плотностью $10^{12}$~м$^{-2}$ ($10^8$~см$^{-2}$):
$r\approx(\pi\rho)^{-1/2}\approx0{,}6$~мкм, $\alpha\approx3$,
$\tau\approx3\cdot(0{,}6\cdot10^{-6})^2/10^{-18}\approx10^6$~с~— около
10~суток.
\item Те же дислокации при $T=400$~К: $D_v$ больше в $\sim10^3$ раз, и
$\tau\approx10^3$~с~— меньше получаса.
\end{itemize}
Видно, как сильно скорость отжига зависит от температуры и от густоты
сети стоков.
\end{primer}

\begin{podrobnee}[Роль источника и состояния поверхности]
Каждый шаг вакансии~— переход через барьер $\Evm$, но её рождение
источником требует ещё работы $A\le\EvF$. Поэтому при нагреве скорость
заполнения объёма вакансиями иногда ограничивает не диффузия, а
\emph{источник}: характерное время рождения
$\tau_{\text{ист}}\approx\nu_0^{-1}\exp\bigl((A+\Evm)/kT\bigr)$ может
превысить время диффузии. У стока, напротив, барьера поглощения либо нет,
либо он мал, и свойства стока менее существенны.

Поглощение вакансий всегда смещает сток (граница или поверхность
«съедает» вакансии и сдвигается), но обычно это несущественно. В сплавах
вакансии могут приносить к стоку примесь и «засорять» его. Для стока на
свободную поверхность важно её состояние: на воздухе металлы покрыты
оксидной плёнкой толщиной 2--5~нм. От её сплошности и сцепления с
металлом зависит, выйдут ли вакансии наружу или будут скапливаться под
плёнкой, образуя поры.
\end{podrobnee}

\subsubsection{Дрейф дефектов в поле сил}

Сток можно «ускорить» дальнодействующим полем. Если энергия частицы $E$
зависит от координаты, то на неё действует сила $\vb F=-\nabla E$, и к
случайному блужданию добавляется направленный \term{дрейф} со скоростью
$\vb v=\mu\vb F=-\mu\nabla E$, где $\mu$~— подвижность. Найдём $\mu$.

\begin{shagi}
\item При концентрации частиц $C$ дрейф создаёт поток $\vb j_1=C\vb v=
-C\mu\nabla E$.
\item В равновесии этот поток уравновешивается встречным диффузионным
потоком $\vb j_2=-D\nabla C$: $\;C\mu\,\dfrac{\d E}{\d x}=-D\dfrac{\d C}{\d x}$.
\item С другой стороны, в равновесии частицы распределены по
Больцману: $C(x)=C(0)\exp\bigl(-E(x)/kT\bigr)$, откуда
$\d C/\d x=-(C/kT)\,\d E/\d x$.
\item Сравнивая, получаем $\mu=D/kT$.
\end{shagi}

\begin{keyeq}\label{eq:einstein}
\mu=\frac{D}{kT}\quad\text{(соотношение Эйнштейна)},\qquad
\vb j=-\frac{CD}{kT}\nabla E .
\end{keyeq}

Для вакансий в поле напряжений энергия $E=\Omega_vp$, где $\Omega_v$~—
объём вакансии, $p$~— гидростатическое давление. Поэтому
\[
\vb j_v=-\frac{C_vD_v\Omega_v}{kT}\nabla p:
\]
вакансии дрейфуют в сторону меньшего давления, т.\,е. в растянутые
области.

Если вакансии непрерывно создаются внешним воздействием (облучением,
пластической деформацией, полиморфным превращением) со скоростью
$\dot C_v$, то устанавливается стационарное пересыщение
$\Delta C_v\approx\tau\dot C_v$. Оно особенно велико при низких
температурах, когда $\tau$ велико.

\subsubsection{Как определить энергию активации по кривым отжига}

На опыте концентрацию дефектов обычно не видят напрямую, а измеряют
какое-нибудь свойство $P$, пропорциональное ей (например, добавочное
электросопротивление, п.~\ref{ssec:exp-methods}). Пусть отжиг определяется
одним процессом с одной энергией активации $E$. Тогда
\begin{keyeq}\label{eq:anneal-kin}
\frac{\d C}{\d t}=-F(C)\,K_0\Ex{-\frac{E}{kT}},
\end{keyeq}
где $F(C)$~— некоторая функция концентрации (вид её зависит от механизма:
например, $F=C$ для ухода на стоки, $F=C^2$ для рекомбинации пар), $K_0$~—
предэкспоненциальный множитель. Важно, что сделаны два допущения:
\begin{enumerate}[label=\asbuk*)]
\item больцмановский множитель $\exp(-E/kT)$ входит в скорость
сомножителем;
\item энергия активации не зависит от концентрации.
\end{enumerate}
Тогда $E$ можно найти, \emph{не зная функцию $F(C)$}. Идея всех методов
одна: сравнить, за какое время (или при какой температуре) достигается
\emph{одна и та же} концентрация в разных условиях. При равных $C$
множитель $F(C)$ одинаков и сокращается.

\paragraph{Метод сечений}
Отжигают серию одинаковых образцов (с одинаковой начальной концентрацией
$C_0$) изотермически при разных температурах $T_1,T_2,\dots$
(рис.~\ref{fig:cut-method},\,а). Интегрируя~\eqref{eq:anneal-kin},
получаем
\[
-\int_{C_0}^{C}\frac{\d C}{F(C)}=K_0\,t\,\Ex{-\frac{E}{kT}} .
\]
Проведём «сечение» $C=C_1=\text{const}$. Левая часть для всех кривых
одинакова, значит, $t_i\exp(-E/kT_i)=\text{const}$, или
\begin{equation}\label{eq:cut}
\ln t_i=\text{const}+\frac{E}{k}\cdot\frac{1}{T_i} .
\end{equation}
Время достижения концентрации $C_1$ в координатах $\ln t$\,--\,$1/T$
ложится на прямую с наклоном $E/k$ (рис.~\ref{fig:cut-method},\,б).
Сечение на другом уровне $C_2$ даёт параллельную прямую~— это проверка
того, что энергия активации действительно одна.

\begin{figure}[ht]
\centering
\begin{subfigure}[b]{0.48\textwidth}\centering
\begin{tikzpicture}
\begin{axis}[width=\linewidth, height=5cm, xmin=0, xmax=10, ymin=0, ymax=1.05,
  xlabel={$t$}, ylabel={$C/C_0$}, xtick=\empty, ytick={0.5}, yticklabels={$C_1$}]
  \addplot[fmRed, thick, domain=0:10, samples=80] {exp(-1.2*x)};
  \addplot[fmAmber, thick, domain=0:10, samples=80] {exp(-0.5*x)};
  \addplot[fmGreen, thick, domain=0:10, samples=80] {exp(-0.22*x)};
  \addplot[fmBlue, thick, domain=0:10, samples=80] {exp(-0.1*x)};
  \draw[dashed] (axis cs:0,0.5) -- (axis cs:10,0.5);
  \addplot[only marks, mark=*, mark size=1.6pt] coordinates {(0.578,0.5) (1.386,0.5) (3.151,0.5) (6.931,0.5)};
  \node[font=\footnotesize, fmRed] at (axis cs:1.0,0.15) {$T_1$};
  \node[font=\footnotesize, fmBlue] at (axis cs:8.6,0.55) {$T_4$};
  \node[font=\footnotesize] at (axis cs:6.5,0.92) {$T_1>T_2>T_3>T_4$};
\end{axis}
\end{tikzpicture}
\caption{}
\end{subfigure}\hfill
\begin{subfigure}[b]{0.48\textwidth}\centering
\begin{tikzpicture}
\begin{axis}[width=\linewidth, height=5cm, xmin=0, xmax=1, ymin=0, ymax=1,
  xlabel={$1/T$}, ylabel={$\ln t$}, xtick=\empty, ytick=\empty]
  \addplot[black, thick, domain=0.05:0.95] {0.1+0.75*x};
  \addplot[black, thick, dashed, domain=0.05:0.95] {0.0+0.75*x+0.2};
  \addplot[only marks, mark=*, mark size=1.8pt] coordinates {(0.15,0.2125) (0.35,0.3625) (0.6,0.55) (0.85,0.7375)};
  \node[font=\footnotesize, rotate=33] at (axis cs:0.45,0.6) {$C=C_2$};
  \node[font=\footnotesize, rotate=33] at (axis cs:0.6,0.47) {$C=C_1$};
  \draw (axis cs:0.6,0.55) -- (axis cs:0.8,0.55) -- (axis cs:0.8,0.7);
  \node[font=\footnotesize, below] at (axis cs:0.75,0.53) {наклон $E/k$};
\end{axis}
\end{tikzpicture}
\caption{}
\end{subfigure}
\caption{Метод сечений: \textit{а}~— изотермические кривые отжига при
разных температурах и сечение $C=C_1$; \textit{б}~— времена достижения
$C_1$ в координатах Аррениуса. Сечения на разных уровнях дают
параллельные прямые}
\label{fig:cut-method}
\end{figure}

\paragraph{Метод отношения наклонов}
Используют один образец. В момент, когда концентрация равна $C^*$,
температуру скачком меняют от $T_1$ до $T_2$ (рис.~\ref{fig:slope-ratio}).
Наклоны кривой непосредственно до и после скачка, согласно
\eqref{eq:anneal-kin}, равны $R_{1,2}=F(C^*)K_0\exp(-E/kT_{1,2})$. Их
отношение не содержит $F(C^*)$:
\begin{equation}\label{eq:slope-ratio}
\ln\frac{R_1}{R_2}=-\frac{E}{k}\left(\frac1{T_1}-\frac1{T_2}\right) .
\end{equation}
Метод даёт одинаковое $E$ при разных $C^*$, если процесс действительно
один. На практике касательные трудно провести точно, поэтому опыт
повторяют несколько раз.

\begin{figure}[ht]
\centering
\begin{tikzpicture}
\begin{axis}[width=0.55\textwidth, height=4.8cm, xmin=0, xmax=10, ymin=0, ymax=1.05,
  xlabel={$t$}, ylabel={$C$}, xtick={4}, xticklabels={$t^*$}, ytick={0.67}, yticklabels={$C^*$}]
  \addplot[fmBlue, very thick, domain=0:4, samples=40] {exp(-0.1*x)};
  \addplot[fmRed, very thick, domain=4:10, samples=40] {exp(-0.4)*exp(-0.35*(x-4))};
  \addplot[fmBlue, dashed, domain=2.5:6.5] {0.6703-0.067*(x-4)};
  \addplot[fmRed, dashed, domain=3.2:5.0] {0.6703-0.2346*(x-4)};
  \draw[dotted] (axis cs:4,0) -- (axis cs:4,0.67);
  \node[fmBlue, font=\footnotesize] at (axis cs:1.6,0.97) {$T_1$};
  \node[fmRed, font=\footnotesize] at (axis cs:7.2,0.25) {$T_2>T_1$};
  \node[fmBlue, font=\footnotesize] at (axis cs:7.3,0.66) {наклон $R_1$};
  \node[fmRed, font=\footnotesize, anchor=west] at (axis cs:4.9,0.4) {наклон $R_2$};
\end{axis}
\end{tikzpicture}
\caption{Метод отношения наклонов: при $C=C^*$ температуру скачком
повышают от $T_1$ до $T_2$ и сравнивают наклоны кривой до и после
скачка}
\label{fig:slope-ratio}
\end{figure}

\paragraph{Нагрев с постоянной скоростью}
Образец с дефектами нагревают с постоянной скоростью $q$: $T=qt$, и
непрерывно измеряют свойство $P$. Пока температура низкая, скорость
отжига ничтожна; затем она резко растёт, и в узком интервале около
характерной температуры $T_a$ исчезают практически все дефекты данной
стадии; после этого $\d P/\d t$ снова падает до нуля. Поэтому производная
$\d P/\d T$ имеет вид пика (рис.~\ref{fig:const-rate}).

\begin{figure}[ht]
\centering
\begin{tikzpicture}
\begin{axis}[name=left, width=0.42\textwidth, height=4.6cm, xmin=200, xmax=400, ymin=0, ymax=1.1,
  xlabel={$T$}, ylabel={$P$}, xtick={300}, xticklabels={$T_a$}, ytick=\empty]
  \addplot[fmBlue, very thick, domain=200:400, samples=120]
    {exp(-1e12*x*x/9284*exp(-9284/x))};
\end{axis}
\begin{axis}[at={(left.east)}, anchor=west, xshift=9mm, width=0.42\textwidth, height=4.6cm,
  xmin=200, xmax=400, ymin=0, xlabel={$T$}, ylabel={$-\d P/\d T$}, xtick={300}, xticklabels={$T_a$}, ytick=\empty]
  \addplot[fmRed, very thick, domain=200:400, samples=200]
    {1e12*exp(-9284/x)*exp(-1e12*x*x/9284*exp(-9284/x))};
\end{axis}
\end{tikzpicture}
\caption{Отжиг при нагреве с постоянной скоростью: слева~— изменение
свойства $P$, справа~— его производная по температуре. Дефекты данной
стадии исчезают в узком интервале около $T_a$}
\label{fig:const-rate}
\end{figure}

Если нагрев ступенчатый~— образец выдерживают одинаковое время при каждой
из последовательно повышающихся температур,~— такой отжиг называют
\term{изохронным}. При малых ступенях он эквивалентен непрерывному
нагреву.

Энергию активации находят по двум опытам с разными скоростями нагрева
$q_1$ и $q_2$. Переходя в~\eqref{eq:anneal-kin} от времени к температуре
($\d t=\d T/q$) и интегрируя (при $E\gg kT$), получаем
\[
\int_{C_0}^{C}\frac{\d C}{F(C)}\approx-\frac{K_0}{q}\cdot\frac{kT^2}{E}\Ex{-\frac{E}{kT}} .
\]
Одинаковая концентрация $C$ достигается в первом опыте при $T_1$, во
втором при $T_2$; левые части равны, и
\begin{equation}\label{eq:const-rate}
\ln\!\left(\frac{q_2T_1^2}{q_1T_2^2}\right)=\frac{E}{k}\left(\frac1{T_1}-\frac1{T_2}\right) .
\end{equation}

\begin{podrobnee}[Сочетание изохронного и изотермического отжига]
Пусть образец отжигают изохронно: выдержки одинаковой длительности
$\Delta t$ при температурах $T_1<T_2<\dots$ Введём «приведённое время»
$\theta=t\exp(-E/kT)$: согласно~\eqref{eq:anneal-kin}, концентрация $C$
зависит только от $\theta$. За $i$-ю выдержку приведённое время
увеличивается на $\Delta\theta_i=\Delta t\,\exp(-E/kT_i)$, а концентрация
меняется от $C_{i-1}$ до $C_i$.

Возьмём теперь изотермическую кривую отжига при температуре $T_a$ и
найдём на ней время $\Delta\tau_i$, за которое концентрация меняется
от того же $C_{i-1}$ до того же $C_i$. Поскольку $C$ зависит только от
$\theta$, это же приращение $\Delta\theta_i=\Delta\tau_i\exp(-E/kT_a)$.
Приравнивая, получаем
\[
\ln\frac{\Delta\tau_i}{\Delta t}=\frac{E}{k}\left(\frac{1}{T_a}-\frac{1}{T_i}\right).
\]
В координатах $\ln\Delta\tau_i$\,--\,$1/T_i$ точки ложатся на прямую с
наклоном $-E/k$, если процесс характеризуется одной энергией активации.
Если одновременно идут два процесса с разными энергиями, этот (как и любой
из описанных) метод даёт «кажущуюся» энергию, зависящую от $C$; в
областях, где преобладает один процесс, она от $C$ не зависит.
\end{podrobnee}

\begin{itogi}
\begin{itemize}
\item Пары Френкеля (вакансия + межузельный атом) рождаются при облучении;
тепловые вакансии~— по механизму Шоттки на поверхностях, границах,
дислокациях. Те же места в пересыщенном кристалле работают как стоки.
\item Отжиг~— уход избыточных дефектов на стоки или их рекомбинация.
\item Пересыщение спадает экспоненциально с $\tau=\alpha r^2/D_v\propto
\exp(\Evm/kT)$; доминирует ближайший сток.
\item Подвижность в поле сил $\mu=D/kT$ (Эйнштейн); вакансии дрейфуют в
растянутые области.
\item Энергию активации отжига находят методами сечений, отношения
наклонов, разных скоростей нагрева~— сравнивая условия достижения
одинаковой концентрации.
\end{itemize}
\end{itogi}

\begin{proverka}
\item Почему механизм Френкеля не работает в тепловом равновесии?
\item Как изменится время отжига вакансий, если плотность дислокаций
увеличить в 100 раз?
\item Почему в методе сечений не нужно знать функцию $F(C)$?
\item Что означает непараллельность прямых $\ln t$\,--\,$1/T$ для разных
сечений?
\end{proverka}
